\documentclass[a4paper,12pt]{letter}
\usepackage{graphicx}
\usepackage[left=0.3in, right=0.3in, top=0.0in, bottom=0.2in]{geometry}

\begin{document}

\begin{letter}{Editor-in-Chief \\
Journal of Environmental Radioactivity \\
Elsevier}

\opening{Dear Editor-in-Chief,}

We submit for your consideration an original research article titled
\textbf{``Low-Methoxyl Pectin--Ca$^{2+}$ Gelation as a Sample
Solidification Protocol for Marine Radioactivity Monitoring by HPGe
Gamma Spectrometry''} for publication in the
\textit{Journal of Environmental Radioactivity}.

The motivation for this work came directly from a practical gap in
shipboard radionuclide monitoring.
Since the Fukushima Daiichi accident, the demand for rapid, on-site
seawater gamma spectrometry during post-accident surveys and routine
naval environmental monitoring has grown substantially.
The standard approach---placing a liquid seawater container on an HPGe
detector aboard a vessel---breaks down because vessel motion causes
the sample surface to oscillate continuously, introducing unpredictable
counting efficiency errors that cannot easily be corrected in the field.
Liquid samples also risk spillage in the confined spaces of a shipboard
instrument room.

Our solution is to convert seawater into a mechanically stable gel
before measurement, using low-methoxyl (LM) pectin and a controlled
addition of CaCl$_2\cdot$2H$_2$O.
The gelation proceeds at room temperature using simple reagents and
completes within 2~h.
The resulting fixed geometry is intended to reduce geometry variation
and leakage risk on mobile platforms.

We validated this approach against Korea Institute of Nuclear Safety
(KINS) certified proficiency-test reference standards across three
monitoring rounds (2019, 2022, 2024), covering four radionuclides:
$^{137}$Cs, $^{134}$Cs, $^{241}$Am, and $^{60}$Co across a
gamma-energy range of 59.5--1173.2~keV.
Mean paired recovery rates were 99.9\% for $^{137}$Cs, 88.8\% for
$^{134}$Cs, 93.7\% for $^{241}$Am, and 94.8\% for $^{60}$Co.
Sixteen of 17 paired measurements remained within 80--120\% recovery.
The manuscript transparently reports the single low $^{134}$Cs paired
result and separates solidification recovery from bias relative to the
KINS certified values.

We believe this work is a good fit for the \textit{Journal of
Environmental Radioactivity} because it addresses a direct obstacle
to shipboard marine radioactivity monitoring, provides quantitative
performance data validated against certified reference standards,
and offers a simple protocol accessible to monitoring groups without
specialised chemistry infrastructure.
To our knowledge, this is the first quantitative evaluation of
pectin-gel solidification of seawater for HPGe gamma spectrometry.

We confirm that this manuscript has not been published previously
and is not currently under consideration elsewhere.
All authors have approved the manuscript and agree with its submission.
The authors declare no conflicts of interest.

We look forward to your consideration.

\closing{Sincerely,}
\vspace{-1.5cm}
\includegraphics[height=1.5cm]{signature.png}\\[0.5em]
\textbf{Ku Kang, Ph.D.\ (Corresponding Author)} \\
CBRN Defense Research Institute \\
Seoul 06796, Republic of Korea \\
Email: bisu9082@gmail.com

\end{letter}

\end{document}
